\begin{resumo}[Abstract]
	\begin{otherlanguage*}{english}
		This work presents the development of a chess simulator integrated with the Stockfish artificial intelligence engine, using the Godot Engine. The main objective was to build a functional, customizable, and highly competitive application, solving the technical challenge of integrating a complex artificial intelligence into a modern game engine. The adopted methodology involved creating the graphical interface and the rule validation logic, such as castling and \textit{en passant}, using the native GDScript language. To process the virtual opponent's moves, the Python language was used as \textit{middleware}, responsible for establishing asynchronous communication between the Godot frontend and the Stockfish engine, operating through the \textit{Universal Chess Interface} (UCI) protocol. Additionally, a dynamic visual theme system based on dictionaries and data arrays was implemented, allowing for real-time customization of the graphical interface. The results obtained demonstrate that the proposed architecture ensures fluid gameplay without compromising computational performance during the calculation of the best moves. It is concluded that the integration between Godot, Python, and Stockfish forms a robust, efficient, and highly modular foundation for the development of digital board games of high technical complexity.

		\vspace*{1\onelineskip}
 
		\textbf{Keywords}: Chess. Godot Engine. Python. Artificial Intelligence. Stockfish.
	\end{otherlanguage*}
\end{resumo}